\documentclass[11pt]{article}

\usepackage[utf8]{inputenc}
\usepackage[T1]{fontenc}
\usepackage{geometry}
\usepackage{amsmath}
\usepackage{amssymb}
\usepackage{booktabs}
\usepackage{hyperref}
\usepackage{xurl}
\usepackage{xcolor}
\usepackage{listings}
\usepackage{enumitem}
\usepackage{parskip}
\usepackage{graphicx}
\graphicspath{{../images/}}

\geometry{margin=1in}

\lstset{
  language=Java,
  basicstyle=\ttfamily\small,
  keywordstyle=\color{blue},
  commentstyle=\color{gray},
  stringstyle=\color{teal},
  showstringspaces=false,
  columns=fullflexible,
  frame=single,
  framerule=0.2pt,
  breaklines=true
}

\setlist[itemize]{topsep=0.3em,itemsep=0.2em}
\setlist[enumerate]{topsep=0.3em,itemsep=0.2em}

% Simple per-section source line.
\newcommand{\notesources}[1]{\par\noindent{\small\color{gray}\textit{Sources:} \sloppy #1\par}}

% Unit-wise source strings for this subject (used by slides/notes).
% Path is relative to the lecture `latex/` folder (the compilation CWD).
% OOPS with Java (CSEG1044) — unit-wise sources
%
% These are short, student-facing reference strings intended to be used with:
%   - \slidesources{...}  (in beamer slides)
%   - \notesources{...}   (in lecture notes)
%
% Style inspiration: other subjects in My-Subjects/ use semicolon-separated sources
% and include an "accessed" date for web documentation.

\newcommand{\OOPSUnitISources}{Schildt, \textit{Java: A Beginner's Guide} (10e), McGraw-Hill (Java fundamentals + OOP basics).; Oracle Java Tutorials: Getting Started + Language Basics + Classes and Objects (accessed \today).; Java SE API docs (e.g., \texttt{String}, \texttt{Scanner}) (accessed \today).}

\newcommand{\OOPSUnitIISources}{Schildt, \textit{Java: The Complete Reference} (12e), McGraw-Hill (classes, inheritance, packages, interfaces).; Bloch, \textit{Effective Java} (3e), Addison-Wesley, 2018 (design choices; interfaces vs abstract classes).; Oracle Java Tutorials: Classes and Objects + Interfaces + Packages (accessed \today).; Java SE API docs (e.g., \texttt{Comparable}, \texttt{Runnable}) (accessed \today).}

\newcommand{\OOPSUnitIIISources}{Schildt, \textit{Java: The Complete Reference} (12e) (nested/inner classes, exceptions, threads, I/O).; Oracle Java Tutorials: Nested Classes + Exceptions + Concurrency + Basic I/O (accessed \today).; Java SE API docs (e.g., \texttt{Thread}, \texttt{AutoCloseable}) (accessed \today).}

\newcommand{\OOPSUnitIVSources}{Schildt, \textit{Java: The Complete Reference} (12e) (generics, lambdas, Swing, JDBC).; Bloch, \textit{Effective Java} (3e) (generics + lambdas).; Oracle Java Tutorials: Generics + Lambda Expressions + Swing + JDBC Basics (accessed \today).; Java SE API docs (e.g., \texttt{java.util.function}, \texttt{javax.swing}, \texttt{java.sql}) (accessed \today).}

\newcommand{\OOPSUnitVSources}{Schildt, \textit{Java: The Complete Reference} (12e) (Collections Framework + wrapper classes).; Oracle Java docs: Collections Framework overview (accessed \today).; Java SE API docs (\texttt{java.util} collections; wrapper classes in \texttt{java.lang}) (accessed \today).}

\newcommand{\OOPSUnitVISources}{Git SCM: \textit{Pro Git} (2e) (version control workflow), accessed \today.; GitHub Docs (repositories, issues, pull requests), accessed \today.; Oracle Java Tutorials: Swing + JDBC (accessed \today).; JUnit 5 User Guide (accessed \today).; Apache Maven Project: Getting Started (accessed \today).; Gradle User Manual: Getting Started (accessed \today).; UML class diagrams: UML 2.x overview/spec (accessed \today).}

\newcommand{\OOPSMiscWhyInterfacesSources}{Schildt, \textit{Java: The Complete Reference} (12e) (interfaces).; Bloch, \textit{Effective Java} (3e) (prefer interfaces where appropriate).; Oracle Java Tutorials: Interfaces (accessed \today).; Java SE API docs (e.g., \texttt{Comparable}, \texttt{Runnable}) (accessed \today).}



\title{Object Oriented Programming (CSEG1044)\\Unit 06 -- Lecture 03 Notes\\UML Class Design + Package Structure}
\author{Tofik Ali}
\date{\today}

\begin{document}
\maketitle
\tableofcontents

\section{Lecture Overview}
Now that you have use cases, you must design the class structure of the system.
This lecture focuses on:
\begin{itemize}
  \item identifying classes and responsibilities,
  \item connecting classes with correct relationships,
  \item reducing coupling using interfaces/abstractions,
  \item organizing code into packages.
\end{itemize}
\notesources{\OOPSUnitVISources}

\section{Syllabus Alignment (Unit 06)}
\textbf{Unit 06:} UML class diagrams; package/module structure; design for maintainability

\section{Learning Outcomes}
After this lecture, you should be able to:
\begin{itemize}
  \item Convert use cases into a draft UML class diagram (entities + services + UI + DAO).
  \item Identify relationships (has-a, is-a) and responsibilities (avoid God classes).
  \item Propose a clean package structure for your project.
  \item Complete Deliverable-2: class diagram (draft) + package structure.
\end{itemize}

\section{Key Terms (Quick)}
\begin{itemize}
  \item \textbf{Entity/model}: represents real-world data (Resource, Booking, User).
  \item \textbf{Service}: business rules (booking validation, overlap checks).
  \item \textbf{DAO}: database access layer (CRUD + SQL).
  \item \textbf{UI layer}: Swing screens/panels.
  \item \textbf{Association}: two classes related (Booking uses Resource).
  \item \textbf{Composition}: strong has-a (Booking has Resource reference).
  \item \textbf{Inheritance}: is-a relationship (use carefully).
  \item \textbf{Coupling}: how strongly modules depend on each other.
\end{itemize}

\section{Detailed Notes}
\subsection{UML class diagram notation (quick cheat sheet)}
A UML class diagram is a \textbf{communication tool}. In this course, we use it to:
\begin{itemize}
  \item show important classes and their responsibilities,
  \item show relationships between classes (especially in your domain model),
  \item justify separation of UI/service/DAO layers.
\end{itemize}

\textbf{Common notations you should recognize:}
\begin{itemize}
  \item A class is drawn as a box with compartments: \textbf{Name}, \textbf{Attributes}, \textbf{Methods}.
  \item Visibility markers: \texttt{+} public, \texttt{-} private, \texttt{\#} protected.
  \item Association: a line between classes (``uses/has-a'').
  \item Multiplicity: \texttt{1}, \texttt{0..1}, \texttt{*}, \texttt{1..*} (how many objects).
  \item Inheritance (is-a): arrow with a hollow triangle pointing to parent.
  \item Composition (strong has-a): filled diamond at the ``owner'' side (use only when lifecycle is tightly linked).
\end{itemize}
\textbf{Tip:} do not try to include every getter/setter in the diagram. Focus on the design.

\subsection{Step-by-step: find classes from use cases}
A practical method:
\begin{enumerate}
  \item List core use cases (from Deliverable-1).
  \item Extract nouns $\rightarrow$ candidate entities (Resource, Booking, User).
  \item Extract verbs $\rightarrow$ operations (create booking, cancel booking).
  \item Group operations into services (BookingService).
  \item Identify persistence needs $\rightarrow$ DAOs (BookingDao, ResourceDao).
  \item Identify UI screens/panels (AddResourcePanel, BookingPanel, ListBookingsPanel).
\end{enumerate}

\subsection{Responsibilities (avoid God class)}
Each class should have one main responsibility.
Bad sign:
\begin{itemize}
  \item One class does UI + SQL + validation + printing.
\end{itemize}
Better separation:
\begin{itemize}
  \item UI reads input and calls service.
  \item Service enforces rules and calls DAO.
  \item DAO only performs DB operations.
\end{itemize}

\subsection{Relationships (has-a vs is-a)}
\begin{itemize}
  \item Prefer composition (has-a) for most domain relationships.
  \item Use inheritance only for true is-a (rare in small projects).
  \item Interfaces are useful to reduce coupling (e.g., DAO interface).
\end{itemize}

\subsection{Interfaces/abstractions to reduce coupling}
Examples:
\begin{itemize}
  \item \texttt{ResourceRepository} interface with methods \texttt{save/findAll}.
  \item Implementation \texttt{JdbcResourceRepository} uses JDBC.
  \item UI/service depends on interface, not implementation (easy to replace/mock).
\end{itemize}

\subsection{Worked mapping: ``Create Booking'' to classes/methods}
To make the class diagram practical, map one use case to the methods you will implement:
\begin{itemize}
  \item \textbf{UI layer:} \texttt{CreateBookingPanel.onSaveClicked()}
  \item \textbf{Service layer:} \texttt{BookingService.createBooking(...)}\par
  validations: mandatory fields; start $<$ end; overlap rule.
  \item \textbf{DAO layer:} \texttt{BookingDao.insert(...)} and \texttt{BookingDao.findOverlaps(...)}
  \item \textbf{Model:} \texttt{Booking} object carries data between layers
\end{itemize}
\textbf{Outcome:} your class diagram should clearly show where each responsibility lives.

\subsection{Package structure (recommended)}
Keep a simple layered package structure:
\begin{itemize}
  \item \texttt{com.upes.crm.model} --- entities
  \item \texttt{com.upes.crm.dao} --- JDBC DAOs
  \item \texttt{com.upes.crm.service} --- business rules
  \item \texttt{com.upes.crm.ui} --- Swing UI
  \item \texttt{com.upes.crm.util} --- helpers (DB connection, validation)
\end{itemize}
\notesources{\OOPSUnitVISources}

\section{Mini Example (Campus resource booking class list)}
\textbf{Entities:}
\begin{itemize}
  \item \texttt{Resource(id, name, type, capacity)}
  \item \texttt{Booking(id, resourceId, userId, startTime, endTime, status)}
  \item \texttt{User(id, name, role)}
\end{itemize}

\textbf{Services:}
\begin{itemize}
  \item \texttt{BookingService}: overlap checks; validate time range; call DAO.
\end{itemize}

\textbf{DAOs:}
\begin{itemize}
  \item \texttt{ResourceDao}, \texttt{BookingDao}, \texttt{UserDao}
\end{itemize}

\textbf{UI:}
\begin{itemize}
  \item \texttt{MainFrame}, \texttt{AddResourcePanel}, \texttt{CreateBookingPanel}, \texttt{ListBookingsPanel}
\end{itemize}

\section{Deliverable-2 checklist}
\begin{itemize}
  \item UML class diagram draft (10--20 classes is typical for this project size)
  \item Show key relationships (especially Booking $\rightarrow$ Resource)
  \item Package structure documented (folder tree or package list)
  \item Updated issues: one issue per major class/module group (optional)
\end{itemize}

\section{In-Class Activity (evidence)}
\begin{itemize}
  \item Upload class diagram and package structure to repo.
  \item Create PR if you add initial project skeleton folders/packages.
  \item Assign issues to teammates based on modules (UI/DB/service).
\end{itemize}

\section{In-Class Exercises (with brief solutions)}
\begin{enumerate}
  \item From your use cases, list 5 nouns (candidate classes).\par
  \textbf{Sample:} Resource, Booking, User, Room, InventoryItem (depends on domain).

  \item For one use case, list 3 methods that must exist.\par
  \textbf{Sample:} \texttt{createBooking}, \texttt{findOverlaps}, \texttt{insert}.

  \item Which layer should contain overlap validation: UI, service, or DAO?\par
  \textbf{Answer:} service layer (business rule), with DAO providing data to check overlaps.
\end{enumerate}

\section{Self-Study Practice (no solutions)}
\begin{enumerate}
  \item Draw a draft UML class diagram for your project (10--20 classes is typical).
  \item Write a package structure for your code (model/dao/service/ui/util).
  \item Choose 2 use cases and map them to UI/service/DAO methods (like the worked mapping).
\end{enumerate}

\section{Checkpoint Questions (with sample answers)}
\textbf{Q1.} Why separate UI, service, and DAO layers?\par
\textbf{Sample answer:} Separation reduces coupling and makes code easier to maintain and test. UI changes should not require SQL changes, and DB changes should not break UI logic directly.

\vspace{0.6em}
\textbf{Q2.} What is the benefit of interfaces in this design?\par
\textbf{Sample answer:} Interfaces let you program to a contract, so you can swap implementations (e.g., JDBC vs in-memory) and reduce dependency on concrete classes.

\end{document}
